%!TEX root = ../thesis.tex
%*******************************************************************************
%*********************************** First Chapter *****************************
%*******************************************************************************

\chapter{Organisms as utility maximisers}

\ifpdf
    \graphicspath{{Chapter1/Figs/Raster/}{Chapter1/Figs/PDF/}{Chapter1/Figs/}}
\else
    \graphicspath{{Chapter1/Figs/Vector/}{Chapter1/Figs/}}
\fi


%********************************** %First Section  **************************************
\section{An Introduction to Utility}

The choice of all possible actions in the world is potentially infinite, yet each organism only has a limited amount of time, effort, and (in the case of humans) money they can expend. As such, a key function of the brain is to accumulate the evidence needed to whittle down the choices and eventually decide on the single course of action that will benefit the organism most over the duration of its remaining lifespan. On the simplest level, organisms are compelled to travel or grow toward increasing gradients of the chemicals that they require to survive and propagate their genes. In evolutionary biological terms, they maximise their ‘reproductive fitness’.

In \textit{Caenorhabditis elegans}, the decision to move forward or backward is simply an integration of evidence for or against a chemical in its vicinity (\cite{tanimotoCalciumDynamicsRegulating2017}). Such decisions are perceptual in that they deal with some external, observable quantity in the outside world. Regardless of whether the worm makes the decision to move away (in the case of this study) from a noxious chemical, there exists a unidimensional scale of the decision variable (i.e., the strength of concentration gradient) against which behaviour can be regressed against.

Rhesus macaques and human participants in neuropsychological studies can produce similar perceptual decisions when presented with an array of randomly moving dots and the chance to win a reward (juice or money) if they can guess the correct aggregate direction (\cite{shadlenNeuralBasisPerceptual2001}). A competent programmer can define exactly the coherence of the movement of dots without having to know anything about the organism participating in the task\footnote{The examples used here are mostly to make a rhetorical point to introduce the idea of external unitary scales of decisions, and the lack thereof in economic decision-making. The point could well be made that all perceptual decisions are a subset of economic decisions. For a brief discussion on economic decision-making as psychophysics, see Appendix A and \cite{lucePossiblePsychophysicalLaws1959}}.

In economic tasks, there is no such external dimension which dictates a correct choice between alternatives and the difficulty of that choice. Even in token economies developed by humans, people frequently make choices which do not maximise monetary value. For example, it is not uncommon for individuals to both seek risk and buy lottery tickets and avoid risk by paying insurance premiums, despite both having a negative expected value in units of currency (\cite{friedmanUtilityAnalysisChoices1948}). Even for certain goods, it is plain to see that a cheeseburger has more benefit to a starving man than a piece of jewellery which may be many times its value\footnote{Of course, this only considers value-in-use, rather value-in-exchange for the purposes of a simplified, illustrative example. We also skip over the idea that it is the marginal utility rather than the utility per se that causes the starving man to value food highly for the purposes of our point. See \cite{whiteDoctoringAdamSmith2002} for a critical take on this common economic fable that goes beyond the scope of this thesis.}. Instead of comparing items against cost, people instead factor the monetary cost into some utility function, which they use to maximise their own internal subjective value for goods.

This tendency for individuals to seek actions that maximise a utility function, rather than value, was first described by Daniel Bernoulli in the $18^{\text{th}}$ century as an answer to the St. Petersburg 'paradox' (\cite{bernoulliExpositionNewTheory1954}). In this thought experiment, a player is offered the opportunity to buy into a lottery where a coin will be flipped n times until it lands on heads. At that point, the player will be paid out $2^{n}$ dollars. The first five terms of the expected value (EV\nomenclature{EV}{Expected Value}) of this game are equal to the chance of each outcome ($p(x_{n})$) multiplied by the magnitude of that outcome ($x_{n}$), and are as follows:

\begin{equation}
\begin{aligned}
EV_n &= p(x_n) \cdot x_n \\
EV &= \sum_{n=1}^{\infty} \left(\frac{1}{2}\right)^n \cdot 2^n \\
EV &= \left(\frac{1}{2} \cdot 2\right) + \left(\frac{1}{4} \cdot 4\right) + \left(\frac{1}{8} \cdot 8\right) + \left(\frac{1}{16} \cdot 16\right) + \left(\frac{1}{32} \cdot 32\right) \ldots = 1 + 1 + 1 + 1 + 1 + \ldots
\label{eq:st_petersburg}
\end{aligned}
\end{equation}

For example, if a player bought into the lottery for \$10 and the coin landed on heads on the third flip, the player would win \$8 and make a total loss of \$2. However, since the summation term in Equation~\ref{eq:st_petersburg} is infinitely long, the expected value of the lottery becomes infinite. That is, a ‘rational’ player should be willing to spend an infinite amount to play this game. Clearly, this defies common sense; it seems clear to an observer that any person who gambled away their life savings buying into the game would likely end up disappointed. In reality, Bernoulli remarks that 'any fairly reasonable man would sell his chance, with great pleasure, for twenty ducats’\footnote{A medieval unit of currency. 20 ducats would have been a substantial sum on the order of hundreds/thousands of dollars in today’s currency. A more recent recreation of the St Petersburg Paradox reveals that modern undergraduates place a median value on the game of \$1.50 (\cite{haydenFictiveRewardSignals2009}) though this uses a modified version where a participant wins $\$2^{(n-1)}$ i.e. half what they would win in the version presented here.}. Instead of judging the expected value of the lottery, he supposed that ‘people with common sense evaluate money in proportion to the utility they can obtain from it’.

This formulation of utility stipulates that the subjective value of possible wins is different from person to person. For instance, though a \$10 bill lying on the sidewalk is a net profit for any person who comes across it, it is likely to elicit much greater joy in a poor man than one who is already rich. If the expected value of the St Petersburg lottery is given in Equation~\ref{eq:st_petersburg}, the expected utility (EU\nomenclature{EU}{Expected Utility}) for a given individual $i$ is given in Equations~\ref{eq:expected_utility1} and \ref{eq:expected_utility2}:

\begin{align}
EU_{i,n} &= p(x_n) \cdot v_i(x_n)
\label{eq:expected_utility1}
\end{align}

\begin{align}
EU_i &= \sum_{n=1}^{\infty} \left(\frac{1}{2}\right)^n \cdot v_i(2^n)
\label{eq:expected_utility2}
\end{align}

Where $\upsilon$ is a function that transforms the objective payment $x_{n}$ to a subjective utility for an individual $i$ studied. Bernoulli’s insight to ‘solve’ the St Petersburg paradox was to imagine a concave utility function: that as an individual’s wealth increases, the marginal utility of winning an additional $x_{n}$ dollars decreases until it eventually plateaus\footnote{In modern terminology the concavity of the utility function forms Gossen’s First Law, also known as the ‘law of diminishing marginal utility’, described by the $19^{\text{th}}$ century economist (\cite{gossen1983laws}).}. A rich individual who could afford to pay \$10,000 would gain comparatively little by increasing their wealth by, e.g. \$6,000 (winning on the $14^{\text{th}}$ flip), which has a proportionally tiny chance of occurring. Bernoulli imagines this utility function to be logarithmic, i.e., that expected utility is proportional to the log of the expected value; however, as the exact shape of the utility curve is unique, the utility of the St Petersburg lottery (or any good, for that matter) varies from individual to individual.

In the nineteenth century, economists dreamed of directly measuring utility using a ‘hedonimeter’ (\cite{edgeworth1881mathematical}) and developed tasks to study the revealed preference of individuals. By presenting subjects with 'bundles' composed of amounts of two goods A and B and studying the choices made between, they were able to measure ordinal utility. For instance, in Figure \ref{fig:indifference}C, a subject is indifferent, that is, has no preference, between a bundle $X$ (which consists of 4 units of good $A$ and 1 unit of good $B$) and a bundle $Y$ (which comprises 1 unit of good $A$ and 5 units of good $B$). If such a choice were presented an infinite number of times to this individual, we would expect them to choose at random and pick each bundle an equal number of times. Non-intersecting indifference curves can be built, representing a monotonic increase in utility as preferential and indifferent choices are made.

\begin{figure}[p] 
\centering    
\includegraphics[width=0.98\textwidth]{Chapter1/Redone_Figs/utility_curves.png}
\caption[Observations of an organism’s binary choices between bundles of goods allows for the construction of utility curves under revealed preference theory.]{\textbf{Observations of an organism’s binary choices between bundles of goods allows for the construction of utility curves under revealed preference theory.} In A-C) an organism chooses between bundles of goods (apples $A$ and bananas $B$). In A) the offered bundle of 4$A$ and 1$B$ is more desirable than an offer of 2$A$ and 1$B$, and so the first bundle is chosen. In B) an offer of 5$B$ and a single $A$ is dominant over an offer of 3$B$ and a 1$A$ and so the former is chosen. In C) the two offers (4$A$ and 1$B$ vs. 5$B$ and a 1$A$) are both preferred equally and will be chosen between as if at random. The organism is \textit{indifferent} between these two choices and they lie on the same utility curve. In D) an example indifference map is generated from the choices in 1A-C). Each point along a curve is a potential choice which carries identical utility to the organism. E.g., the choice offered in C) is highlighted on the lowest curve. As curves are drawn away from the origin, the utility they represent to the organism increases; the points along a curve offer more utility to the organism than any possible point below the curve. Each distinct curve is referred to as an ‘indifference curve’. All choices in this figure are hypothetical. Adapted from \cite{pastorMonkeysChooseIf2017}.}
\label{fig:indifference}
\end{figure}

To generate these ‘indifference curves’, only two axioms are required. The first axiom of completeness is described by:

\begin{equation}
\begin{aligned}
\forall X, Y: X \succ Y, Y \succ X, \text{ or } X \sim Y
\label{eq:axiom_indifference}
\end{aligned}
\end{equation}

For any bundles $X$ and $Y$, either a bundle $X$ is preferred to a bundle $Y$, $Y$ is preferred to $X$, or bundles $X$ and $Y$ are valued equally. The second axiom is that of transitivity, namely:

\begin{equation}
\begin{aligned}
\text{if } X \succ Y \text{ and } Y \succ Z \Rightarrow X \succ Z
\label{eq:axiom_transitivity}
\end{aligned}
\end{equation}

If a bundle $X$ is preferred to a bundle $Y$, and a bundle $Y$ is preferred to a bundle $Z$, then bundle $X$ is preferred to bundle $Z$.\footnote{Stated here only for strict preference for brevity; the same relation holds more generally for weak preference, $X \succeq Y \text{ and } Y \succeq Z \Rightarrow X \succeq Z$, which subsumes the strict and indifferent cases above.} A third implication is that we ignore goods that might have disutility and assume that a greater magnitude of a good(s) in a bundle is always preferential (as illustrated in Figure \ref{fig:indifference}B).

However, this ordinal revealed preference theory of utility can only explain choices made between certain (riskless) choices. It was not until 1944 that von Neumann and Morgenstern showed that a cardinal utility function could be sampled from a subject's choices between risky gambles in their seminal work (\cite{vonNeumann1944}). The authors proposed four axioms (a further two plus those above in Equations \ref{eq:axiom_indifference} and \ref{eq:axiom_transitivity}) which any decision maker must satisfy if they can be said to be behaving as if they are maximising their expected utility (as defined in Equations \ref{eq:expected_utility1} and \ref{eq:expected_utility2}\footnote{Sometimes von Neumann-Morgenstern expected utility is differentiated from this equation as
$$EU_{i,n} = p(x_n) \cdot U_i(x_n)$$
For the purposes of this thesis, we consider the differences in notation unimportant.}). These final two axioms are the axioms of continuity and independence, and are shown in Equations \ref{eq:axiom_continuity} and \ref{eq:axiom_independence}:

\begin{equation}
\begin{aligned}
\forall X \succ Y \succ Z, \exists p \in [0,1]: pX + (1-p)Z \sim Y
\label{eq:axiom_continuity}
\end{aligned}
\end{equation}

For all choices where $X$ is preferred to $Y$, which itself is preferred to $Z$ (and so, by transitivity, $X$ is preferred to $Z$), there must exist some probability $p$ for which the composite of $p \cdot X$ and $(1-p) \cdot Z$ which is equal in preference to the second option $Y$.The continuity axiom dictates that the ordered choices of revealed preference theory can be compared on a continuous scale of utility where probability and magnitude of outcome are integrated.

\begin{equation}
\begin{aligned}
\forall X \succeq Y, \Rightarrow \forall Z \text{ and } \forall p \in [0,1]: pX + (1-p)Z \succeq pY + (1-p)Z
\label{eq:axiom_independence}
\end{aligned}
\end{equation}

In all cases where $X$ is at least as preferable as $Y$, then for any possible third option $Z$ and any possible probability $p$, the bundling of either the probability $p$ or the third option $Z$ to both sides of the original choice should have no effect on the preference of the first option over the second. If the first three axioms of this Expected Utility Theory (completeness, transitivity, and continuity) describe the shape of a utility function, the independence axiom states that preferences should be static. The order choices are presented in, or the addition of decoys to the presented choice, should have no effect on which is chosen.

Any economic decision-maker that satisfies these four axioms can be said to be making choices as if they are trying to maximise their utility. Thus, the formalisation by von Neumann and Morgenstern allowed for the measurement of an organism's utility from the binary choices\footnote{Throughout this thesis we refer to experimental paradigms where a subject indicates a choice between two goods, or bundles, as a ‘binary choice task’. Three such binary choices are illustrated in Figure \ref{fig:indifference}A-C). This is compared in Chapter 2 to an auction task in which a subject indicates a value they would be willing to pay for a good/bundle in terms of another other good.} they make between bundles of goods with given probabilities, and underpins the last half of a century of research into both human (\cite{mostellerExperimentalMeasurementUtility1951}) and animal (\cite{ferraritonioloReliablePopulationCode2023}; \cite{genestUtilityFunctionsPredict2016}; \cite{pastorMonkeysChooseIf2017}) studies of economic decision-making. Many violations of the four axioms have been discovered during this time, most notably the paradoxes of \cite{allaisComportementLHommeRationnel1953} and \cite{ellsbergRiskAmbiguitySavage1961}; \cite{machinaChoiceUncertaintyProblems1987} and Appendix A provide broader accounts of these limitations and the development of more descriptive models over the past half century. Nonetheless, Expected Utility Theory was seminal to research into economic decision-making and shaped many of the experimental paradigms employed in the literature.

In modern behavioural economics and judgment and decision-making research, however, Expected Utility Theory has largely been superseded by frameworks such as Cumulative Prospect Theory \cite{tverskyAdvancesProspectTheory1992}, which has two key aspects that better account for documented violations. First, probabilities are treated subjectively: individuals systematically overweight low-probability outcomes and underweight high-probability ones, producing an S-shaped distortion of the probability scale \cite{kahnemanProspectTheoryAnalysis1979}. Second, utility is evaluated relative to a reference point rather than in absolute terms, and the disutility of a loss exceeds the utility of an equivalent gain, a phenomenon termed \textit{loss aversion}. This asymmetry produces a characteristic kink in the value function at the reference point, where the function is steeper in the loss domain than the gain domain. A fuller treatment, including the Allais paradox which motivated Prospect Theory, is given in Appendix A (Section \ref{sec:prospect}).

\clearpage
\section{Hedonimetry and the Brain}\label{subsec:hedonimetry}

If behavioural science tells us that animals act as if they are maximising the utility they receive, it is not unreasonable to suggest that a function of the brain is to track the utility of offered choices. If such a collection (or collections) of cells that encoded a common currency of utility across choice domains could be identified, and their activity deciphered, it might prove possible to realise Edgeworth’s dream of a hedonimeter to fully measure the utility function of any decision. With complete measurement and knowledge of the utility signal in the brain, the fantasy of behavioural science of predicting the choices of any subject from any set of options offered would be achieved.

The first evidence that such a population might exist came from studies of rats that would repeatedly self-stimulate via a lever connected to an electrode implanted in the basal forebrain (\cite{oldsPositiveReinforcementProduced1954}). There is nothing intrinsic to any physical object that is ‘rewarding’; the concept of an object’s rewarding value comes from the reinforcement it produces to collect/do/eat more of it (\cite{salamoneEffortrelatedFunctionsNucleus2007}). These implanted rats, which will indefatigably press a lever to continue self-stimulating, show the extreme positive reinforcement generated by stimulation of the basal forebrain, implying that activity of cells in the vicinity are conveying reward.

Later studies of intracranial self-stimulation (ICS\nomenclature{ICS}{Intracranial Self-Stimulation}) revealed the ipsilateral pharmacological oxidopamine (6-OHDA) lesioning of the ascending dopaminergic fibres reduced the compulsion\footnote{  Especially in the ICS literature it is difficult to anthropomorphise the rats as receiving ‘pleasure’ from the self-stimulation. Even when studying physiological markers which correlate with pleasure in humans, it is not clear to us that any model organism necessarily has a positive emotional valence associated with their choices/behaviour in neuroeconomic studies. Throughout this chapter, we use words such as ‘pleasure’ and ‘reward’ while detailing the progression of neuroscience into economic decision-making, building up to studies which instead correlate neuronal activity with the more formal, neutral term ‘utility’. Throughout the rest of the thesis, we ascribe an animal’s motivations to maximising this neutral quantity of utility. For a brief review of ‘wanting’ vs ‘liking’ motivations in organisms making behavioural choices, see Appendix B.} of rats to self-stimulate (\cite{fibigerRoleDopamineIntracranial1987}, \cite{phillipsDopaminergicSubstratesIntracranial1976}). This implication of midbrain dopaminergic neurons which project to the striatum and frontal cortices in the reinforcing activity of ICS was further strengthened by studies which showed that levels of ICS were correlated with the presence of dopamine metabolites (\cite{fibigerRoleDopamineIntracranial1987}) and that animals were most reinforced to self-stimulate when electrode tips were placed within the ventral tegmental area (VTA\nomenclature{VTA}{Ventral Tegmental Area}) or substantia nigra pars compacta (SNc\nomenclature{SNc}{Substantia Nigra pars compacta}) with the volume of ICS correlating with the density of dopaminergic neuron cell bodies at this point (\cite{corbettIntracranialSelfstimulationRelation1980}).

In addition to self-stimulating for current, rats will also press a lever to administer drugs of abuse intravenously. Injection of pimozide, a dopamine receptor antagonist, greatly increases the self-administration of amphetamines by animals in an effort to compensate for the reduced efficacy per volume of drug (\cite{yokelIncreasedLeverPressing1975}). Similarly, 6-OHDA lesions to dopaminergic projections to the nucleus accumbens in the ventral striatum have also been shown to extinguish the self-administration of amphetamines and cocaine in rats (\cite{lynessDestructionDopaminergicNerve1979}; \cite{robertsExtinctionRecoveryCocaine1980}) further implying dopamine’s role in the reinforcement of ‘rewarding’ behaviours. The incentive-sensitisation theory builds from these findings to posit that the mesolimbic dopamine pathway plays a central role in generating the craving for addictive drugs (\cite{robinsonNeuralBasisDrug1993}); maladaptive or not, the dopaminergic pathways robustly generate motivations which lead to choices, and thus have an important role in any measurement of inferred utility.

Physiologically, disorders of dopamine in humans have a marked effect on a patient’s motivations. Parkinson’s disease, which involves the destruction of dopaminergic neurons in the substantia nigra, is well known to induce apathy (a motivational deficit) in sufferers (\cite{denbrokApathyParkinsonsDisease2015}), which can be lifted in many patients by the application of the dopamine precursor L-Dopa (\cite{czerneckiMotivationRewardParkinsons2002}). In many patients, attempts to address the reduction in dopaminergic activity can lead to the opposite outcome of compulsive behaviour such as excessive gambling (\cite{grossetProblematicGamblingDopamine2006}) and hypersexuality (\cite{romitoTransientManiaHypersexuality2002}). Dopaminergic dysfunction is also implicated in schizophrenia, a condition in which sufferers also regularly suffer from apathy. In such patients, the ventral striatal responses to reward signals are dampened (\cite{juckelDysfunctionVentralStriatal2006}).

As mentioned above, choices and actions are only ‘rewarding’ insofar as they can reinforce the said behaviour. Therefore, there is a crucial link between reward and learning. One influential model of this associative learning is the Rescorla-Wagner model (\cite{rescorla1972theory}), which describes how the strength of association between a conditioned stimulus (CS) and an unconditioned stimulus (US) is updated on a trial-by-trial basis through their co-occurrence. The model is associative rather than causal: it does not encode cause-effect relationships between stimuli, nor does it make predictions about future states. Changes in associative strength are computed within each trial only. The update in associative strength between a CS and US can be written as:

\begin{equation}
\begin{aligned}
\Delta V_x^n &= \alpha \beta \left( \lambda - \sum V_{all} \right) \\
V_x^{n+1} &= V_x^n + \Delta V_x^n \\
x &\in all
\end{aligned}
\label{eq:rescorla}
\end{equation}

Where $V_x^n$ represents the strength of association between a CS-US pairing $x$ at trial $n$. This is updated for trial $n+1$ by taking the strength of association at time n and adding the change ($\Delta$) in strength of association that occurs during trial n. This change is determined by the salience of the conditioned stimulus in $x$ (a red sky in the morning might predict future rain as well as the distant rumble of thunder and lightning, but it is much less noticeable), specified by a parameter $\alpha$, and the learning rate of the subject for the unconditioned stimulus $\beta$. These parameters then scale the prediction error term inside the brackets; the difference between the maximum associative strength $\lambda$ and the combined sum of associative strength for all stimuli present $\sum V_{all}$ . I.e., the change in associative strength on trial $n$ is proportional to the difference between what is expected, $\lambda$, and what actually happens, $\sum V_{all}$.

The success of the Rescorla-Wagner model in predicting empirical learning findings such as overshadowing (the stronger conditioning to the most salient group of conditioned stimuli) and blocking (the deficit in conditioning to a second conditioned stimulus when paired with an already perfectly conditioned stimulus; \cite{millerAssessmentRescorlaWagnerModel1995}).

This model was the foundation of later 'temporal difference' (TD) models that made two important contributions relevant to the neuroscientific study of utility (\cite{ludvigEvaluatingTDModel2012}; \cite{suttonModernTheoryAdaptive1981}). First, the Rescorla-Wagner model treats the trial as the unit of time, which is unsuitable for physiological learning which must happen continuously in real-time. It is conceivable that something learned in the first half of a trial might be vital in the later half, or even on some later trial. Secondly, TD models employ a Markov approach to updating predictions where the difference in predictions at each time step are compared to each other and then to the unconditioned stimulus, also referred to as ‘second order conditioning’ (\cite{suttonTimeDerivativeModelsPavlovian1990}). For an intuitive example, consider an American Football team moving the ball down the pitch in separate plays (time-step) where the end result of each play carries a different future prediction of the unconditioned reward points that the team might score (\cite{dawValueLearningReinforcement2014}).

Temporal difference learning employs the Bellman equation (\cite{bellmanDynamicProgrammingStatistical1957}) to keep calculating the total value of a state $x$ at time $t$. Any such state has a reward immediately available $r_t$ (a pass into the endzone, in American football terms), but this may be vanishingly small (/unlikely). Any change in the state over time opens up new rewards to be sampled and new states to progress to. The reward $r_{t+1}$ that a team can expect to sample will be predicated upon their current state, and so on as the decision tree unfolds. The Bellman equation can be written as:

\begin{equation}
\begin{aligned}
V(x_t) &= r_t + \gamma \mathbb{E}[V(x_{t+1}) | x_t] \\
\text{or equivalently} \\
V(x_t) &= \mathbb{E}[r_t + \gamma r_{t+1} + \gamma^2 r_{t+2} + \gamma^3 r_{t+3} + \ldots | x_t]
\end{aligned}
\label{eq:bellman}
\end{equation}

Any such huge decision tree in Equation \ref{eq:bellman} can be simplified down into smaller steps that iterate recursively. The temporal difference error signal can then be expressed as the difference between the predictions of any agent of their current state and time and the former state and time.

These theoretical models had a profound effect on the search for a neural substrate of reward when reconciled with electrophysiological recordings of non-human primate midbrain dopaminergic neurons (\cite{schultzNeuralSubstratePrediction1997}). Such dopaminergic neurons typically have stereotyped responses to unpredicted rewards, showing phasic activity with a duration of less than 200ms and a latency of <100ms (\cite{ljungbergResponsesMonkeyDopamine1992}). This activity scales with the magnitude (typically measured in ml of a pleasurable sugary juice) of reward, and with the reward's unpredictability (\cite{bayerMidbrainDopamineNeurons2005}); these neurons show the greatest increase in stimulus-related firing rate\footnote{  For a colloquial explanation of the signal: dopamine neurons typically have a tonic firing rate of roughly 5Hz (imagine a fast bass drum beat). For a view of what 'firing' means, see Figure \ref{fig:isi_waveforms} showing the average trace of a 'spike' from a representative dopaminergic neuron we recorded in this project. Upon delivery of an unexpected reward, midbrain dopamine neurons produce a very brief 'drumroll' of the same spikes with a much shorter latency between. The described increase in firing rate occurs within this 'drumroll' before the neuron returns to its 'resting' tonic frequency} when the organism receives the most surprising, most preferred kind (e.g. flavour), and largest reward. Such responses of a single dopaminergic neuron to an unpredictable reward of juice recorded from the midbrain of a rhesus macaque over multiple trials are shown in Figure \ref{fig:da_raster}A).

Crucially, this 'quantitative encoding' (\cite{bayerMidbrainDopamineNeurons2005}) of the desirability extends to cues which predict a future reward (\cite{mirenowiczImportanceUnpredictabilityReward1994}). Indeed, as a conditioned stimulus for reward becomes established, dopaminergic neurons gradually increase the stimulus-related phasic activity corresponding to the presentation of the conditioned stimuli and decrease stimulus-related phasic activity at the time of the reward delivery. For a perfectly temporally predictive cue-reward relationship, the only change in dopaminergic firing rate is seen when the conditioned stimulus is presented and there is no change in the reward-related dopaminergic firing rate (\cite{schultzNeuralSubstratePrediction1997}). Changes in the firing rate of a single dopaminergic neuron evoked by a fully conditioned stimulus and reward are shown in Figure \ref{fig:da_raster}B).

This change in the firing rates for these two stimuli presentations can be explained using the prediction error paradigm described above. A subject to whom it has been signalled that a reward will occur with 100\% reliability in exactly $n$ seconds will have no prediction error at the actual receipt of the reward. This temporal-difference 'reward-prediction error' (RPE\nomenclature{RPE}{Reward Prediction Error}) is what dopaminergic neurons signal with their change in activity (\cite{dardenneBOLDResponsesReflecting2008}; \cite{schultzNeuralSubstratePrediction1997}). The compatibility of this signal with temporal-difference (rather than the foundational Rescorla-Wagner) models is supported by a study by \cite{montagueFrameworkMesencephalicDopamine1996}, who showed that rewards could be predicted from sequences of conditioned stimuli. Such sequences are only compatible with the higher-order prediction errors of temporal-difference learning.

\begin{figure}[p] 
\centering    
\includesvg[width=0.85\textwidth,inkscapelatex=false]{Chapter1/Redone_Figs/RPE_CS_remade.svg}
\caption[Single dopaminergic cell responses to rewards and conditioned stimuli.]{\textbf{Single dopaminergic cell responses to rewards and conditioned stimuli.} A), the responses of a single non-human primate dopamine cell we recorded to an unexpected juice reward (R) are shown. Reward was delivered at time 0 and the responses of the cell are shown as a peri-stimulus time histogram (PSTH\nomenclature{PSTH}{Peri-Stimulus Time Histogram}; top) and as a raster plot (bottom) of single spikes from this cell. The phasic response of the cell can be seen starting at 100ms and continuing for 200ms. As our task was complex, panels B) and C) show simulated data to illustrate dopamine neurons activity upon simple conditioning of a stimulus (B) and when a predicted reward is not forthcoming (C). For biological data, see \cite{schultzNeuralSubstratePrediction1997}. B) a conditioned stimulus (CS) cue is presented at time point 0, predicting a reward (R) at 2500ms. The firing rate of the simulated dopaminergic neuron increases at time 0 with the CS presentation rather than at the delivery of the reward. C) the same CS cue is shown at time 0, but the predicted reward is never delivered (R-). Simulated neurons maintain phasic activity at time 0 but shows a decrease in firing rate upon the omission of the predicted reward.}
\label{fig:da_raster}
\end{figure}

If the change in firing rate of these neurons is correlated to the reward-prediction error of a pair of conditioned and unconditioned stimuli, it should hold for when this error is negative, i.e. when a reward is predicted, but does not arrive. Such a relationship is observed (\cite{schultzNeuralSubstratePrediction1997}) and an illustrative example is presented in Figure \ref{fig:da_raster}C). In addition, punishments and conditioned stimuli for future punishments have also been shown to inhibit the dopaminergic firing rate (\cite{mirenowiczImportanceUnpredictabilityReward1994})\footnote{'Aversive' stimuli can create an initial positive increase in dopaminergic firing rate thought to be mediated by a 'saliency' signal (\cite{schultzDopamineRewardPredictionerror2016}). This saliency signalling of dopaminergic occurs earlier (within the first 120ms of stimulus presentation) and is correlated with the perceptual intensity of the stimulus (e.g. brightness, loudness), the novelty of the stimulus (though novelty alone is not sufficient to generate it; \cite{toblerAdaptiveCodingReward2005}), and whether it shares characteristics (e.g. modality) of other conditioned stimuli for rewards (\cite{dayAssociativeLearningMediates2007}).}.

\clearpage
\section{A Dopaminergic Neural Substrate of Utility}

That the change in dopaminergic neuron firing rate encodes a reward prediction error, with responses correlated with the physical components of expected value, i.e., the magnitude and probability of the reward is a well established finding (\cite{fiorilloDiscreteCodingReward2003}; \cite{toblerAdaptiveCodingReward2005}). However, as discussed, physical components alone cannot account for the level of utility received by an organism. Instead, expected utility theory formalises the calculation of the utility present in any choice by observing the behaviours of a choosing subject.

Recent behavioural experiments in mice (\cite{kimOptogeneticMimicryTransient2012}), rats (\cite{changBriefOptogeneticInhibition2016}; \cite{steinbergCausalLinkPrediction2013}), and monkeys (\cite{staufferDopamineNeuronSpecificOptogenetic2016a}) have shown that optical stimulation of 'mimic' dopaminergic RPE activity is sufficient to drive learning and choice behaviour. In contrast, the disruption of the ability of dopaminergic neurons to produce burst firing by genetic inactivation of NMDA\nomenclature{NMDA}{N-Methyl-D-Aspartate (receptor)} receptors led to mutant mice that had attenuated performance in a bandit choice task (\cite{zweifelDisruptionNMDARdependentBurst2009}), implicating these cells as a neuronal substrate of utility. More directly still, dopamine transients have been shown to be both necessary and sufficient for the acquisition of associative learning, including learning that goes beyond simple model-free associations to model-based ones (\cite{sharpeDopamineTransientsAre2017}).

\cite{lakDopaminePredictionError2014} furthered this by studying the dopaminergic activity of midbrain neurons in rhesus macaques making choices that differed by magnitude, probability, and type of consumable good (either flavour of juice or juice versus mashed banana). Monkeys showed complete and transitive choices, i.e. choosing as if satisfying some integrated quantity (utility) rather than separate reward dimensions. The RPE activity of individual neurons recorded in this study reflected this unitary scale, rather than the separate attributes of the goods, suggesting that the prediction error of dopaminergic cells is, in fact, a utility prediction error.

In addition to physical components, another important dimension in which rewards can vary is the effort and time required to receive the good, that is, 'a bird in the hand is worth two in the bush.' If dopamine signalling incorporates all the dimensions of a possible reward on a unitary scale, these too should be encoded. \cite{dayAssociativeLearningMediates2007} used voltammetry to study dopamine released by midbrain projections to the nucleus accumbens in rats. In forced choice trials, significantly less dopamine was released (with the implication that therefore there is also less neuronal spiking activity from the upstream dopaminergic neurons, but see\footnote{A simplification. Dopamine release is not a simple linear function of dopaminergic cell firing rate (\cite{mohebiDissociableDopamineDynamics2019}). In fact, dopamine release may not even fully necessitate anything upstream of dopaminergic varicosities, as cholinergic stimulation in the striatum seems to be able to shunt dopamine release (\cite{mohebiDopamineReleaseDrives2020}).}) when the obligated good was delayed or required repeated lever presses (at a high effort cost). Single cell dopaminergic activity in macaques has since been shown to be reduced hyperbolically as rewards are discounted over time (\cite{kobayashiInfluenceRewardDelays2008}), and a recent study has extended the idea of folding effort to the unitary scale of utility encoding in macaque dopaminergic neurons (\cite{burrellWorthWorkMonkeys2023}).

Finally, employing an experimental paradigm using gambles to construct the utility function (the 'fractal method'; \cite{machinaChoiceUncertaintyProblems1987}), \cite{staufferDopamineRewardPrediction2014} found that the normalized dopaminergic phasic activity associated with the reward increased non-linearly as the expected value of a gamble increased linearly. Such nonlinearities in utility versus value are the foundational theoretical finding of Bernoulli's work on the St. Petersburg paradox. That the increase in dopaminergic cell firing mirrors this asymmetry is more evidence that the activity of these dopaminergic projections signal a utility-prediction error.

We focus on the utility-prediction error account of dopamine throughout this thesis, but it is important to note that this is not the only current interpretation of the computational role of dopaminergic activity. A number of findings in the wider literature sit uneasily with a strict reading of this theory. Most notably, dopamine neurons can respond strongly to outcomes that are already clearly and reliably predicted, which a fully-trained, canonical prediction error signal should not do (\cite{howeProlongedDopamineSignalling2013}; \cite{hamidMesolimbicDopamineSignals2016}). These findings have motivated several alternative or extended accounts of dopaminergic computation. One prominent recent proposal holds that dopamine instead signals how reliably one event actually causes or predicts another (working backwards from an outcome to figure out what was responsible for it) rather than a moment-to-moment prediction error over value (\cite{jeongMesolimbicDopamineRelease2022}; \cite{namboodiriWhyDopamineCausal2024}). Others have proposed that dopaminergic prediction errors reflect the animal's own uncertain, evolving belief about what situation or stage of a task it is in, rather than a simple, fixed link between a cue and its outcome (\cite{starkweatherDopamineRewardPrediction2017}). More broadly still, dopaminergic activity is increasingly understood to be far more varied than accounts that put forward just one single interpretation of what it does, with different dopamine neurons shown to respond quite differently depending on the type of task, stimuli, and conditions involved (\cite{dabneyDistributionalCodeValue2020}; \cite{Lowet_Zheng_Meng_Matias_Drugowitsch_Uchida_2025}; \cite{sousaMultidimensionalDistributionalMap2025}). We discuss these alternative accounts in the context of our own results in Section~\ref{subsec:what_does_da_compute}.

Although we have focused on the dopaminergic midbrain in this introduction (and throughout the rest of this thesis), other areas of the brain have also been proposed as encoding economic value. Given the fundamental importance of economic decision-making for the survival and well-being of humans and other animals, it is hardly surprising that these decisions might involve the integration of multiple regions of the brain (\cite{grabenhorstPredictionEconomicChoice2012}; \cite{padoaschioppaNeuronsOrbitofrontalCortex2006}). For more detail on the anatomy of the brain circuits that are believed to underlie reward processing, and evidence for other regions of the brain's involvement in these circuits, see Appendix C.

As a putative neural substrate for utility in the brain, we spend the remainder of this thesis discussing and analysing the behaviour of non-human primates on a task that elicits utility from choices on a trial by trial basis. This behavioural utility is then used to investigate the activity of recorded dopaminergic neurons in the midbrain of these animals.
